%% ============================================================================
%%  applications/template.tex  —  SKELETON RESUME (single source template)
%% ----------------------------------------------------------------------------
%%  This is the ONLY resume template. The orchestrator copies this file into a
%%  position folder, then the writer fills the Experience / Projects / Skills
%%  sections from the context.md bullet pool. It replaces the old per-track
%%  templates (full-stack / ai-ml / dev-ops) — track is now a doctrine-routing
%%  label, not a file choice, so any track (including robotics/autonomy) fills
%%  this same skeleton.
%%
%%  WHAT IS FIXED (do not change): header (\name, \contact) and Education.
%%  WHAT GETS FILLED: Experience entries, Projects entries, Technical Skills.
%%
%%  rSectionEntry signature — 5 positional params, ALWAYS supply all 5 braces:
%%     {1 name}{2 dates}{3 tagline/role/link}{4 technologies}{5 location}
%%     - param 4 (technologies) renders after the name with a " | " separator;
%%       leave it {} for experiences, fill it for projects (the tech-stack line).
%%     - any unused param is an empty {}.
%%  rSet signature — {1 bucket label}{2 comma-separated items}.
%%
%%  PATH NOTE: this template sits at resume/applications/template.tex, beside
%%  template.cls, so the class path is just `template`. On copy into
%%  applications/{year}/{company}/{role}/, the orchestrator rewrites it to
%%  ../../../template (three levels up to applications/, where template.cls lives).
%% ============================================================================

\documentclass[
	11pt
]{../../../template}

\name{Ankur Desai}
\contact{(248) $\hspace{-0.1em}$ 657 $\hspace{-0.4em}$ -- $\hspace{-0.4em}$ 3805 \\ ardusa05@gmail.com \\ \href{https://ardusa.dev/}{www.ardusa.dev} \\ \href{https://linkedin.com/in/ardusa/}{linkedin.com/in/ardusa}}

\begin{document}

	% ---- EDUCATION (fixed; writer does not modify) --------------------------
	\begin{rSection}{E}{ducation}
		\begin{rSectionEntry}
			{Michigan State University}
			{Expected May 2028}
			{B.S. in Computer Science}
			{}
			{East Lansing, MI}

			\item \textbf{GPA:} 3.5 / 4.0
			\item \textbf{Coursework:} Data Structures \& Algorithms, OOP, Computer Architecture, Linear Algebra, Discrete Math
		\end{rSectionEntry}
	\end{rSection}

	% ---- EXPERIENCE (writer fills from context.md pool) ---------------------
	\begin{rSection}{E}{xperience}
		\begin{rSectionEntry}
			{Claude Builder Club @ MSU}
			{Oct 2025 -- Present}
			{Co-Founder \& Vice President}
			{}
			{East Lansing, MI}

			\item Built a hands-off attendance-enforcement pipeline that runs a daily in-database cron job, resolves absentees, and posts targeted Slack reminders, automating a three-strikes rule that previously required manual tracking.
			\item Built the club's mentorship lineage as an interactive 3D radial graph in react-three-fiber, hand-writing a proportional-arc layout that prevents edge crossings and a custom camera that autofits each family to the viewport.
		\end{rSectionEntry}

		\begin{rSectionEntry}
			{MindMosaic}
			{Aug 2025 -- May 2026}
			{Software Engineering Intern}
			{}
			{East Lansing, MI}

			\item Migrated hot read paths off recursive SQL self-joins onto a Neo4j graph layer so relationships stay first-class, cutting query latency 48\% (850ms to 440ms) and making multi-hop traversals 3x faster than the equivalent SQL.
			\item Designed a dual-database layer that keeps transactional records in PostgreSQL while serving relationship-heavy reads from Neo4j, routing each query class to the store that answers it fastest instead of forcing one engine to do both jobs.
			\item Built and shipped an interactive knowledge-graph platform that keeps dense, interconnected content navigable as it grows, delivering a production MVP to real users on a polyglot Java Spring Boot, Neo4j, and PostgreSQL stack.
			\item Owned the Spring Boot backend serving the platform's core CRUD and graph-traversal APIs, structuring endpoints around the dual-store model so the client never has to know which database backs a given read.
		\end{rSectionEntry}

	\end{rSection}

	% ---- PROJECTS (writer fills from context.md pool) -----------------------
	\begin{rSection}{P}{rojects}
		\begin{rSectionEntry}
			{fliks}
			{\href{https://github.com/fliks-gg}{github.com/fliks-gg}}
			{Multi-game platform that crowd-validates gameplay skill and mints shareable certification cards}
			{Go, PostgreSQL, AWS EC2}
			{}

			\item Architected an ML rating service that grades each clip independently and shows its verdict beside the crowd's, with those crowdsourced ratings serving as the labeled data the classifier trains on.
			\item Coordinated transcode jobs through a Postgres-backed queue using SELECT FOR UPDATE SKIP LOCKED instead of bolting on Redis or SQS, keeping the whole system on one datastore and one self-provisioned EC2 box.
			\item Made ratings immutable and gated commenting behind a submitted rating, so the certification signal cannot be retroactively gamed and social participation feeds skill data instead of diluting it.
		\end{rSectionEntry}

		\begin{rSectionEntry}
			{Dadei}
			{\href{https://dadei.app}{dadei.app}}
			{Ambient voice assistant that turns overheard conversation into human-approved calendar, email, and task actions}
			{Python, pgvector}
			{}

			\item Replaced single-signal context lookup with hybrid retrieval over pgvector that fuses vector similarity, lexical overlap, recency decay, and participant overlap, grounding each proposal under a fixed token budget.
			\item Engineered a no-merge-first speaker-ID model using an EMA voice centroid plus a bounded prototype bank with floor-and-margin gating, abstaining into a new identity on ambiguity, attributing 96\% of utterances to the correct speaker.
		\end{rSectionEntry}

		\begin{rSectionEntry}
			{WizViz}
			{\href{https://devpost.com/software/wizviz}{devpost.com/software/wizviz}}
			{Webcam fighting game driven by real-time body-gesture recognition, winner of the Interactive Media track at SpartaHack X}
			{Python, MediaPipe, OpenCV}
			{}

			\item Ran pose inference asynchronously off the render loop with a monotonic timestamp guard that drops out-of-order results, sustaining 60 FPS with sub-20ms input lag instead of blocking on per-frame detection.
			\item Classified five distinct combat gestures from body landmarks with a scale-invariant geometric model normalized by torso length, making detection work across players and camera distances without any labeled training data.
		\end{rSectionEntry}
	\end{rSection}

	% ---- TECHNICAL SKILLS (writer picks 3-5 buckets from context.md) --------
	\begin{rSection}{T}{echnical Skills}
		\begin{rSet}{Languages}{Python, SQL, Java, JavaScript/TypeScript}
		\end{rSet}
		\begin{rSet}{Databases}{PostgreSQL, Neo4j, Redis, pgvector}
		\end{rSet}
		\begin{rSet}{Cloud \& Infrastructure}{AWS (EC2, S3, RDS)}
		\end{rSet}
	\end{rSection}

\end{document}
